\section{The good lower space and multiplicities at a good point}

% BEGIN V86E1
\paragraph{V86E1. Residual congruence selects the original localization.}
At location $i$ let $\sigma$ index a nonzero summand of
$\mathrm{goodPartM}(i)$. Write $\Sigma_0,x_0$ for witnesses of
its defining condition $\mathrm{AtM}$. Apply J86O1 to the upper
theta type $P$, giving an integral character $\psi$ of the common
Hecke algebra $T_c$ for $(j(i),w(i))$. In this case $\psi$ has a
point already at $i$. To see this, take the row list $\Sigma_1$
and augmented base change of $P$ from U86T4. Both $\Sigma_0$
and $\Sigma_1$ localize to the same lower parameter at each split
place. Adding the common string and twisting therefore identifies
the local parameter of $P$ with the augmentation of $\Sigma_0$.
At a tested split place the integral polynomial $p_w$ in
$\mathrm{AtM}$ has this normalized local polynomial as its
complex image. By the coefficient calculation in J86O3, in the
common ring $O_c$ its coefficient of $X^{n_V-b}$ is
\[
    (-1)^b(\mathrm Nw)^{b(b-1)/2}\,\psi(T_{w,b})
      \qquad(1\le b\le n_V).
\]
Here we identify $O_i$ with $O_c$ by the member's
$\mathrm{phO}$, which respects the embeddings into $\mathbb C$.
The reduction of $p_w$ is that of the Hecke polynomial of $x_0$.
Since $w\nmid\ell_i$, its norm is a unit of $O_c$. Cancelling
the displayed unit for each $b$ gives equality of the two residual
characters on the generators. They also agree on constants and
hence on $T_c$. The residual kernel of $\psi$ is precisely
the maximal ideal indexing $i$, by the point description
$\mathrm{Sp.ptX}$ for $x_0$. Thus $\mathrm{Sp.ptX}$ at $i$
itself supplies $y\in X_i$ for $\psi$, since its simultaneous
eigenspace is nonzero by J86O1. The test on all Hecke generators
gives $\mathrm{ptV}(i,P)=y$.

The augmented list of $P$ contains the distinguished string,
so $y\in Z$. The test for $\mathrm{goodSig}(\Sigma_1)$ holds
because $\mathrm{HasParamW}$ identifies the row labels at the
designated split places with the list labels. Thus $y\in Z_{\rm good}$
and the original summand is in $M_{j,\Sigma_1}$ for a list above
that very point. In the opposite direction, fix $x\in Z_{\rm good}$,
any list $\Sigma$ over $x$, and any row with
$\mathrm{HasParamW}(i,\sigma,\Sigma)$. Its finite type is good
by G85M3. The list and $x$ witness $\mathrm{AtM}$: use the
Hecke polynomial of $x$ itself at each test place, since its
coefficient image equals the augmented local polynomial by the
base change of a witness of $\mathrm{HasParamX}$ and
$\mathrm{pt\_loc}$. Taking the sums of these inclusions in
both directions proves $\mathrm{RowsC.goodPart}$.
% END V86E1

% BEGIN V86E2
\paragraph{V86E2. A pointwise multiplicity test for both downward conditions.}
Fix a location $i$, a point $x\in Z_{\rm good}(i)$ and
a list $\Sigma$ over $x$. Set $A=M_{j,\Sigma}$,
$E=\mathrm{eigC}(i,x)$, and write $T$ for the original upward
theta operator. Its adjoint $T^*$ is the downward operator, with
the member's specified metrics. J86O3 gives $TA\subseteq E$,
and L86M2 gives injectivity of $T$. Also $T^*TA\subseteq A$:
on every generating lower type the composite acts by its positive
scalar by L86M4. These finite-dimensional spaces consequently satisfy
\[
 \dim A\le\dim E,\qquad E=TA\oplus\big(E\cap(TA)^\perp\big).
\]
If $\dim E\le\dim A$, equality gives $E=TA$. Then $T^*E\subseteq A$,
and for $z=Ta\in E$ the identity $T^*z=0$ implies
$\langle z,z\rangle=\langle Ta,z\rangle=0$.
Conversely, $T^*E\subseteq A$ and injectivity of $T^*$ on $E$
give $\dim E\le\dim A$. Thus the joint assertions of
$\mathrm{downMem}$ and $\mathrm{typeNZ}$ are equivalent to the
single bound
\[
                  \mathrm{mult}(i,x)\ \le\ \dim M_{j,\Sigma}
\]
for all good $x$ and every list over them. Here the identity
$\mathrm{mult}(i,x)=\dim E$ is part of the original point data.

The part of $T^*E$ detected by the named finite types can be
determined before proving this bound. If a nonzero lower type
$\sigma$ pairs with $E$ after applying $T$, its unique upper
type $P$ from U86T3 lies above $x$, by orthogonality of the upper
types. H86P3 gives equal parameters at every split place for
$P$ and a witness for $\mathrm{HasParamX}(i,x,\Sigma)$.
Untwisting and cancelling the common string in the formula of
U86T1 gives the lower localization $\mathrm{isoLoc}(\Sigma)$.
The positivity and rank clauses of $\mathrm{HasParamW}$ follow
from the augmented base-change list on removing the positive
rank-$t$ string. Hence $\sigma$ is among the summands of $A$.
If $B$ is the span of the other named lower types, this proves
$T^*E\perp B$. Moreover $TA\perp TB$ because the types lifting
from $A$ have point $x$ by J86O3 and those lifting from $B$ do
not. The positive scalar identities on the summands of $A$
give $T^*TA=A$, so also $A\perp B$. Therefore
\[
                T^*E\ \subseteq\ A\oplus (A+B)^\perp .
\]
In particular the lower spanning assertion at this row in L86M5
would give $\mathrm{downMem}$ there. For $\mathrm{typeNZ}$, the
positive scalar of $T^*T$ on the lower side controls $TA$.
The possible additional upper multiplicity is exactly
$\dim E-\dim A$. The native joint downward test requires that
this number vanish at every good point.
% END V86E2
